\section{指数映射与豪斯多夫公式}

关于指数映射的第一批研究属于 E. 施图迪与 F. 恩格尔；恩格尔{[}104 b{]}注意到，指数映射对 \(\mathbf{SL}(2, \mathbf{C})\) 不是满射（例如 \(\left( \begin{array}{cc} -1 & a \\ 0 & -1 \end{array} \right)\) 当 \(a \neq 0\) 时不是一个指数），但对 \(\mathbf{GL}(n, \mathbf{C})\) 是满射，从而对 \(\mathbf{PGL}(n, \mathbf{C})\) 也是（后一性质已由施图迪对 \(n = 2\) 指出）；于是 \(\mathbf{SL}(2, \mathbf{C})\) 与 \(\mathbf{PGL}(2, \mathbf{C})\) 给出了一个例子：两个群局部同构，但从整体观点看却非常不同。恩格尔还证明，在其余的经典群添上位似之后，指数映射是满射；这项工作被毛雷尔、施图迪等人重新拾起并继续，而没有带来实质性的新结果。

1899 年，H. 庞加莱（{[}251{]}，v. III, pp.~169-172 与 173-212）从一个不同的观点着手指数映射的研究。他的论文似乎写得仓促，因为在几处他断言连通群的每个元素都是指数，而在别处又给出相反的例子。他的结果主要涉及伴随群：他证明，这种群 G 的半单元素可以是李代数 \(L(G)\) 中无穷多个元素的指数，而非半单元素则可能不是指数。若 \(\operatorname{ad}(X)\) 没有作为 \(2\pi i\) 的非零倍数的特征值，则 exp 在 X 处是平展的。他还证明，若 U 与 V 在 \(L(G)\) 中描出回路，且 W 由连续性定义使得 \(e^{U}.e^{V}=e^{W}\)，则我们未必回到 W 的初始判定。他使用一个留数公式，它本质上化归为

\[
\Phi (\operatorname{ad} X) = \frac {1}{2 \pi i} \int \frac {\Phi (\xi) d \xi}{\xi - \operatorname{ad} X}
\]

其中 \(\operatorname{ad}(X)\) 是一个非零特征值均为单重的半单元素，\(\Phi\) 是一个收敛半径足够大的整级数，积分取在包围 ad X 的特征值的一条回路上；他还研究了当 X 趋于一个具有多重特征值的变换时会发生什么。

在庞加莱的工作之前不久，对 \(W\)——作为 \(U\) 与 \(V\) 的函数——在公式 \(e^U.e^V = e^W\) 中的表达式的寻求，就已经是坎贝尔{[}46 a and b{]}两篇论文的主题。稍后贝克写道：“\ldots 李的理论以显然的方式提示，乘积 \(e^Ue^V\) 形如 \(e^W\)，其中 \(W\) 是 \(U\) 与 \(V\ldots\) 的一个交错级数”。这个主题的后续工作旨在把这一断言精确化，并给出 \(W\) 的一个显式公式（或一种构造方法）（“豪斯多夫公式”）。在坎贝尔与庞加莱之后，帕斯卡、贝克{[}13{]}与豪斯多夫{[}152 a{]}又回到这个问题；每个人都认为其前人的证明不能令人信服；主要的困难在于“交错”一词的含义：所涉及的是所考虑的特定李代数的元素，还是普遍的“符号”表达式？坎贝尔、庞加莱与贝克在这一点上都没有表达清楚。相反，豪斯多夫的论文是完全精确的；他首先在有限个不定元的形式结合（非交换）级数代数中工作，并把 U、V、W 看作这个代数的元素。他借助一个与其前人类似的微分方程论证证明了 W 的存在性。同一个论证又被他用来证明当把其中的不定元换成有限维李代数的元素时级数的收敛性。正如贝克所注意到的、庞加莱也独立注意到的，这一结果可以用来给出李第三定理的一个证明；它澄清了群与李代数之间的对应，例如就换位子群而言。

1947 年，邓金{[}98 a{]}重新研究这个问题，并从一开始就考虑赋范李代数（有限维或非有限维，建立在 R、C 或一个超度量域上），从而得到了豪斯多夫公式的显式系数。\(^{19}\)

